\exercisesection{}

¶ 1) 设 G 为局部紧群，β 为 G 上的一个左 Haar 测度，A 为 G 的一个 β-可测子集，ν 为 G 上一个非零正测度。假设对每个 \(s \in G\) ，sA 都是局部 ν-可忽略的。证明 A 是局部 β-可忽略的。（化归为 A 相对紧且 ν 有界的情形。证明 ν 与 \(\varphi_{A} \cdot \beta\) 可卷积，且 \(\nu *^{\beta} \varphi_{A} = 0\) ，从而 \(0 = \|\nu * \varphi_{A} \cdot \beta\| = \|\nu\| \cdot \|\varphi_{A} \cdot \beta\|\) 。）在承认连续统假设的前提下证明：若不假设 A 是 β-可测的，这一结果可能不成立。（取 \(G = R^{2}\) ，对 ν 取子群 \(R \times \{0\}\) 上的 Haar 测度，并应用 Ch. V, §8 的 Exer. 7 c)。）

\begin{enumerate}
\def\labelenumi{\arabic{enumi})}
\setcounter{enumi}{1}
\tightlist
\item
  设 H 是赋有离散拓扑的加法群 R。设 G 为局部紧群 \(R \times R \times H\) 。设 \(\alpha\) 是 G 上的一个 Haar 测度，\(\beta\) 是 \(R \times H\) 上的一个 Haar 测度，而 \(\mu = \varepsilon_{0} \otimes \beta \in \mathcal{M}(G)\) 。在 G 上构造一个函数 \(f \geqslant 0\) ，它局部 \(\alpha\) -可忽略（从而 \(\mu\) 与 f 可卷积），但 f 的任何平移都不是 \(\mu\) -可测的。（模仿 Ch. V, §3, Exer. 4 的构造。）
\end{enumerate}

¶ 3) 设 G 为局部紧群。

\begin{enumerate}
\def\labelenumi{\alph{enumi})}
\tightlist
\item
  设 \(\mu\) 是 \(\mathbf{G}\) 上一个非零的有界正测度，满足 \(\mu * \mu = \mu\) 。证明支集 \(\mathbf{S}\) （\(\mu\) 的）是紧的。（取 \(f \in \mathcal{K}_{+}(\mathbf{G})\) 且 \(f \neq 0\) ；在 \(\mathbf{G}\) 上选定一个左 Haar 测度，卷积都将关于它取；我们有 \(\mu * f \in \overline{\mathcal{K}(\mathbf{G})}\) ；设 \(x \in \mathbf{S}\) 使得 \((\mu * f)(x) = \sup_{y \in \mathbf{S}} (\mu * f)(y)\) ；证明
\end{enumerate}

\((\mu * f)(y) = (\mu * f)(x)\) 对每个 \(y \in S\) 成立。

\begin{enumerate}
\def\labelenumi{\alph{enumi})}
\setcounter{enumi}{1}
\tightlist
\item
  证明 S 是 G 的紧子群，且 \(\mu\) 是 S 的规范化 Haar 测度。（利用 a)、GT, III, §4 的 Exer. 21 以及 §3 的 Exer. 5。）
\end{enumerate}

\begin{enumerate}
\def\labelenumi{\arabic{enumi})}
\setcounter{enumi}{3}
\tightlist
\item
  设 \(G\) 为局部紧群。对每个 \(t \in \mathbf{R}_+^*\) ，设 \(\mu_t\) 是 \(G\) 上一个非零的有界正测度。假设映射 \(t \mapsto \mu_t\) 对拓扑 \(\sigma(\mathcal{M}^1(G), \overline{\mathcal{K}(G)})\) 连续，且 \(\mu_{s+t} = \mu_s * \mu_t\) （ \(s, t\) 属于 \(\mathbf{R}_+^*\) ）。
\end{enumerate}

\begin{enumerate}
\def\labelenumi{\alph{enumi})}
\item
  证明存在一个数 \(c \in \mathbf{R}\) 使得 \(\| \mu_t \| = \exp(ct)\) 。（注意 \(t \mapsto \| \mu_t \|\) 是下半连续的，且 \(\| \mu_{s+t} \| = \| \mu_s \| \cdot \| \mu_t \|\) 。利用 Prop. 18。）
\item
  假设 \(c = 0\) 。证明 \(\mu_t\) 当 \(t \to 0\) 时弱收敛到 \(G\) 的某个紧子群的规范化 Haar 测度。（利用 \(\mathcal{M}^1(G)\) 的单位球的弱紧性，并对每个弱聚点 \(\mu\) （\(t \mapsto \mu_t\) 关于 \(R_+^*\) 中 0 的邻域滤子的），证明 \(\mu_t * \mu = \mu_t\) 对每个 \(t\) 成立，然后 \(\mu * \mu = \mu\) ；接着应用 Exer. 3。）
\end{enumerate}

¶ 5) 设 G 为在无穷处可数的局部紧群，β 为 G 上的一个左 Haar 测度，\(a \in R_{+}^{*}\) ，且 \((\nu_{u})_{0 \leqslant u \leqslant a}\) 为 G 上满足下列条件的一族正测度：

\begin{enumerate}
\def\labelenumi{(\roman{enumi})}
\item
  \(\nu_0 = \varepsilon_e\) ；\(\nu_{u + v} = \nu_u * \nu_v\) （若 \(u + v \leqslant a\) ）；\(\nu_u = \stackrel{\vee}{\nu}_u\) ；
\item
  对 \(0 < u \leqslant a\) ，\(\nu_{u} = f_{u} \cdot \beta\) ，其中 \(f_{u} \geqslant 0\) 是下半连续的；
\item
  \(\nu_{u}\) 是 \(u\) 的一个淡连续函数。
\end{enumerate}

\begin{enumerate}
\def\labelenumi{\alph{enumi})}
\tightlist
\item
  设 \(f \in \mathcal{K}_{+}(\mathrm{G})\) 。证明对 \(0 < u \leqslant a\) ，
\end{enumerate}

\[
\int \left(f _ {u / 2} * f\right) (z) ^ {2} d \beta (z) = \int f (x) \left(f _ {u} * f\right) (x) d \beta (x).\tag{1}
\]

\begin{enumerate}
\def\labelenumi{\alph{enumi})}
\setcounter{enumi}{1}
\item
  设 \(f \in \mathcal{K}(\mathbf{G})\) 。证明 \(f_{u/2} * f\) 有 \(\beta\) -可积的平方（对 \(0 < u \leqslant a\) ），且 (1) 仍成立。
\item
  设 \(f \in \mathcal{H}(\mathrm{G})\) 满足 \(\nu_{a} * f = 0\) 。证明 \(f = 0\) 。（有 \(\nu_{a/2^n} * f = 0\) 对每个整数 \(n > 0\) （由 \(b\) )），故由 §3, No.~3 的注 1 得 \(f = 0\) 。）
\item
  设 \(\nu \in \mathcal{C}'(\mathrm{G})\) 满足 \(\nu_{a} * \nu = 0\) 。证明 \(\nu = 0\) 。（正则化 \(\nu\) （用 \(\mathcal{H}(\mathrm{G})\) 中的函数）并应用 \(c\) ）。）
\end{enumerate}

\begin{enumerate}
\def\labelenumi{\arabic{enumi})}
\setcounter{enumi}{5}
\item
  设 G 为局部紧群。证明：若 \(\mathcal{K}(\mathrm{G})\) 对卷积是交换的，则 G 是交换群。（用正则化证明 \(\mathcal{C}'(\mathrm{G})\) 是交换的，并把这一点应用到测度 \(\varepsilon_{s}\) 上，其中 \(s \in G\) 。）
\item
  设 \(G\) 为局部紧群，\(\beta\) 为 \(G\) 上的一个左 Haar 测度。证明：代数 \(L^1(G, \beta)\) 有单位元当且仅当 \(G\) 是离散的。（设 \(G\) 非离散，并取 \(f_0 \in L^1(G, \beta)\) 。存在一个紧邻域 \(V\) （\(e\) 的）使得
\end{enumerate}

\[
\int_ {V} | f _ {0} (x) | d \beta (x) <   1.
\]

设 \(\mathbf{U}\) 是 \(e\) 的一个对称紧邻域，使得 \(\mathbf{U}^2 \subset \mathbf{V}\) 。那么对几乎每个 \(x \in \mathbf{U}\) ，

\[
\left| \left(\varphi_ {\mathrm{U}} * f _ {0}\right) (x) \right| = \int_ {\mathrm{U}} \left| f _ {0} \left(y ^ {- 1} x\right) \right| d \beta (y) \leqslant \int_ {\mathrm{V}} \left| f _ {0} (x) \right| d \beta (x) <   1,
\]

因而 \(f_0\) 不是 \(L^1(G, \beta)\) 的单位元。

\begin{enumerate}
\def\labelenumi{\arabic{enumi})}
\setcounter{enumi}{7}
\tightlist
\item
  设 G 为在无穷处可数的局部紧群，连续地从左边作用在一个波兰局部紧空间 T 中。设 \(\nu\) 是 T 上一个在 G 下拟不变的正测度。设 R 是 T 上一个 \(\nu\) -可测的等价关系，与 G 相容。于是存在 (Ch. VI, §3, No.~4, Prop. 2) 一个波兰局部紧空间 B 和一个 T 到 B 中的 \(\nu\) -可测映射 p，使得 \(R\{x,y\}\) 等价于 \(p(x)=p(y)\) 。设 \(\nu'\) 是 \(\nu\) 在 p 下的伪象测度，并设 \(b\mapsto\lambda_{b}\) （ \(b\in B\) ）是 \(\nu\) 关于 R 的一个分解。证明 \(\lambda_{b}\) 对几乎每个 \(b\in B\) 在 G 下拟不变。
\end{enumerate}

（设 \(\chi\) 是 \(\mathbf{G} \times \mathbf{T}\) 上满足 Ch. VII, §1 的 Exer. 13 的条件的函数。证明对每个 \(s \in \mathbf{G}\) ，存在 B 的一个 \(\nu'\) -可忽略子集 \(\mathbf{N}(s)\) ，使得 \(\chi(s^{-1}, \cdot)\) 是局部 \(\lambda_b\) -可积的（对 \(b \notin \mathbf{N}(s)\) ）；为此，注意对每个 \(\psi \in \mathcal{K}(\mathbf{T})\) ，函数 \(x \mapsto \psi(x)\chi(s^{-1}, x)\) 是 \(\lambda_b\) -可积的，除了 \(b\) 属于某个 \(\nu'\) -可忽略集 \(\mathbf{N}(s, \psi)\) 以外，并利用 Ch. VI, §3, No.~1 的引理 1。令 \(\lambda_{b,s}' = 0\) （若 \(b \in \mathbf{N}(s)\) ），并令 \(\lambda_{b,s}' = \chi(s^{-1}, \cdot) \cdot \lambda_b\) （若 \(b \notin \mathbf{N}(s)\) ）。证明映射 \(b \mapsto \lambda_{b,s}'\) 是 \(\nu'\) -适配的（利用 Ch. VI, §3, No.~1 的引理 3），且 \(\gamma(s)\beta = \int \lambda_{b,s}' d\nu'(b)\) 。另一方面证明 \(\gamma(s)\beta = \int \gamma(s)\lambda_b d\nu'(b)\) ，并由此推出：对每个 \(s \in \mathbf{G}\) ，有 \(\gamma(s)\lambda_b = \chi(s^{-1}, \cdot) \cdot \lambda_b\) 对几乎每个 \(b\) ，从而对几乎每个 \(b\) 有 \(\gamma(s)\lambda_b = \chi(s^{-1}, \cdot) \cdot \lambda_b\) 对几乎每个 \(s\) 。然后利用 §4 的 Prop. 17 的推论 2 推断：对几乎每个 \(b\) ，\(\gamma(s)\lambda_b\) 都与 \(\lambda_b\) 等价（对所有 \(s \in \mathbf{G}\) ）。）

证明：若 \(\nu\) 在 G 下以乘子 \(\chi\) 相对不变，则对几乎每个 b，\(\lambda_{b}\) 以乘子 \(\chi\) 相对不变。

\begin{enumerate}
\def\labelenumi{\arabic{enumi})}
\setcounter{enumi}{8}
\item
  设 \(G\) 为局部紧群，\(\beta\) 为 \(G\) 上的一个左 Haar 测度，\(A\) 与 \(B\) 为两个 \(\beta\) -可积集，满足 \(\beta(A) \leqslant \beta(B)\) 且 \(\beta^{*}(A^{-1}) < +\infty\) 。证明存在
\item
  不相交集合 \(N, K_{1}, K_{2}, \ldots\) 覆盖 \(A\) ，其中 \(N \beta\) -可忽略而诸 \(K_{n}\) 是紧的；
\item
  覆盖 B 的不相交集合 \(\mathbf{N}^{\prime}\) 、\(\mathbf{K}_1^{\prime}\) 、\(\mathbf{K}_2^{\prime}\) 、\ldots{} ，其中诸 \(\mathbf{K}_n^{\prime}\) 是紧的；\\
\item
  元素 \(s_n \in G\) 使得 \(\mathbf{K}_n^{\prime} = s_n\mathbf{K}_n\) 。
\end{enumerate}

（利用 \(\beta (x\mathbf{A}\cap \mathbf{B})\) 连续地依赖于 \(x\) 这一事实以及

\[
\int \beta (x \mathrm{A} \cap \mathrm{B}) d \beta (x) = \beta (\mathrm{A} ^ {- 1}) \beta (\mathrm{B}),
\]

证明：若 \(\beta (\mathbf{A})\neq 0\) ，则存在 \(x\in \mathbf{G}\) 使得 \(\beta (x\mathbf{A}\cap \mathbf{B})\neq 0\) 。）

\begin{enumerate}
\def\labelenumi{\arabic{enumi})}
\setcounter{enumi}{9}
\tightlist
\item
  设 \(\mathbf{G}\) 为局部紧群，\(\beta\) 为 \(\mathbf{G}\) 上的一个左 Haar 测度。对任意两个 \(\beta\) -可积集 \(\mathbf{A}, \mathbf{B}\) ，令 \(\rho(\mathbf{A}, \mathbf{B}) = \beta(\mathbf{A} \cup \mathbf{B}) - \beta(\mathbf{A} \cap \mathbf{B})\) 。
\end{enumerate}

\begin{enumerate}
\def\labelenumi{\alph{enumi})}
\item
  设 A 是一个 \(\beta\) -可积集。证明 \(x \mapsto \rho(xA, A)\) 是连续函数。
\item
  设 U 是 e 的一个邻域。证明存在一个紧集 A 和一个数 \(\varepsilon > 0\) ，使得 \(\rho(x\mathbf{A}, \mathbf{A}) < \varepsilon\) 蕴含 \(x \in U\) 。（取 A 为 e 的一个使 \(A \cdot A^{-1} \subset U\) 的邻域。）
\item
  为使 G 的子集 C 相对紧，必须且只需存在一个 \(\beta\) -可积集 A 和一个数 \(a (0 < a < 2\beta(A))\) ，使得 \(x \in C\) 蕴含 \(\rho(xA, A) \leqslant a\) 。
\end{enumerate}

\begin{enumerate}
\def\labelenumi{\arabic{enumi})}
\setcounter{enumi}{10}
\tightlist
\item
  \begin{enumerate}
  \def\labelenumii{\alph{enumii})}
  \tightlist
  \item
    设 G 为由 e 的一个紧邻域生成的局部紧群。设 \(\varphi\) 是 G 的一个非满射连续自同态，它属于 G 的（双连续）自同构所成的群在紧收敛拓扑下的闭包。那么 \(\lim_{\psi\in\mathcal{G},\psi\to\varphi}\mod\psi=0\) 。（设 K 是生成 G 的 e 的一个紧邻域。设 \(\mu\) 是 G 上的一个左 Haar 测度。若 \(\mu(\varphi(K)) > 0\) ，则 \(\varphi(K) \cdot \varphi(K)^{-1}\) 是 \(e\) 在 G 中的一个邻域，因此 \(\varphi(G)\) 是 G 的开子群；对 \(\psi \in \mathcal{G}\) 充分接近 \(\varphi\) ，有 \(\psi(K) \subset \varphi(G)\) ，从而 \(\psi(G) \neq G\) ，这是荒谬的。于是 \(\mu(\varphi(K)) = 0\) 。当 \(\psi \in \mathcal{G}\) 趋于 \(\varphi\) 时，\(\mu(\psi(K))\) 趋于 0。）
  \end{enumerate}
\end{enumerate}

\begin{enumerate}
\def\labelenumi{\alph{enumi})}
\setcounter{enumi}{1}
\tightlist
\item
  设 G 是一个自由交换群，是直和 \(G_{1} \oplus G_{2} \oplus \cdots\) ，其中每个 \(G_{i}\) 都同构于 Z。把 G 看作离散的。设 \(\varphi\) 是 G 的非满射自同态 \((x_{1}, x_{2}, \ldots) \mapsto (0, x_{1}, x_{2}, \ldots)\) 。那么 \(\varphi\) 是 G 的自同构关于紧收敛拓扑的极限，且 G 的每个自同构的模都等于 1。
\end{enumerate}

\begin{enumerate}
\def\labelenumi{\arabic{enumi})}
\setcounter{enumi}{11}
\tightlist
\item
  对 \(t > 0\) 与 \(x \in \mathbf{R}\) ，令 \(F_t(x) = te^{-\pi t^2 x^2}\) 。设 \(f \in \overline{\mathcal{H}(\mathbf{R})}\) 。证明 \(f * F_t\) 一致地趋于 \(f\) （当 \(t\) 趋于 \(+\infty\) 时）。（证明测度 \(F_t(x) dx\) 满足 §2, No.~7 的引理 4 的条件，其中 \(a = 0\) 。）
\end{enumerate}

¶ 13) 设 G 为局部紧群，连续地从左边作用在一个波兰局部紧空间 T 中。设 β 是 G 上的一个左 Haar 测度，ν 是 T 上一个拟不变正测度，χ(s,x) 是 G × T 上一个 \textgreater0 的函数，满足 Ch. VII, §1 的 Exer. 13 的条件。

对 \(f \in \mathrm{L}^p(\mathrm{T}, \nu)\) （ \(1 \leqslant p < +\infty\) ），令

\[
\left(\gamma_ {\chi , p} (s) f\right) (x) = \chi (s ^ {- 1}, x) ^ {1 / p} f (s ^ {- 1} x).
\]

\begin{enumerate}
\def\labelenumi{\alph{enumi})}
\item
  证明对每个 \(s \in G\) ，\(\pmb{\gamma}_{\chi,p}(s)\) 是 \(L^p(T, \nu)\) 的一个等距自同态。证明映射 \(s \mapsto \pmb{\gamma}_{\chi,p}(s)\) 是 \(G\) 在 \(L^p(T, \nu)\) 中的一个线性表示（按 §2, No.~5 的方式论证）。
\item
  设 \(f \in \mathcal{L}^p(\mathbf{T}, \nu)\) ，并设 \(h \in \mathcal{K}(\mathbf{G})\) 。证明函数
\end{enumerate}

\[
(x, s) \mapsto f (s ^ {- 1} x) h (s) \chi (s ^ {- 1}, x) ^ {1 / p}
\]

是 \(p\) 次方可积的（对 \(\beta \otimes \nu\) ）（先从 \(p = 1\) 的情形开始，并利用 §4, No.~1 的引理 1）。证明：若 \(q\) 是与 \(p\) 共轭的指数，且 \(g \in \mathcal{L}^q(\mathrm{T}, \nu)\) ，则 \(f(s^{-1}x)h(s)g(x)\chi(s^{-1},x)^{1/p}\) 对 \(\beta\otimes\nu\) 可积（把 \(h\) 写成 \(h_1h_2\) 的形式，其中 \(h_1,h_2\) 属于 \(\mathcal{K}(\mathrm{G})\) ）。然后由 Lebesgue-Fubini 定理推出：对 \(f\in\mathrm{L}^{p}(\mathrm{T},\nu)\) 与 \(g\in\mathrm{L}^{q}(\mathrm{T},\nu)\) ，函数

\[
s \mapsto \int g (x) (\gamma_ {\chi , p} (s) f) (x) d \nu (x)
\]

是 \(\beta\) -可测的。

\begin{enumerate}
\def\labelenumi{\alph{enumi})}
\setcounter{enumi}{2}
\item
  试利用 §4, No.~6 的 Prop. 18 的推论 2 以及 Ch. VI, §3, No.~1 的引理 1 证明：表示 \(s \mapsto \gamma_{\chi,p}(s)\) （G 在 \(L^p(T, \nu)\) 中的）对 \(1 \leqslant p < +\infty\) 是连续的。
\item
  此外假设对每个 \(s \in G\) ，函数 \(x \mapsto \chi(s^{-1}, x)\) 有界。对 \(f \in L^{p}(T, \nu)\) ，令
\end{enumerate}

\[
\left(\gamma_ {\chi} (s) f\right) (x) = \chi (s ^ {- 1}, x) f (s ^ {- 1} x).
\]

证明 \(s \mapsto \gamma_{\chi}(s)\) 是 \(G\) 在 \(L^p(T, \nu)\) 中的一个连续表示（\(1 \leqslant p < +\infty\) ）。（如同 §2, No.~5 的命题 9 的证明中那样，先证明 \(s \mapsto \gamma_{\chi}(s)\) 是 \(G\) 的一个表示，以 \(L^p(T, \nu)\) 的自同态为值。然后注意：若 \(f \in L^p(T, \nu)\) 且 \(h \in \mathcal{K}(G)\) ，则 §4, No.~1 的引理 1 表明 \(h(s)f(s^{-1}x)\chi(s^{-1}, x)\) 是 \(p\) 次方可积的（对 \(\beta \otimes \nu\) ）。证明的结尾与 c) 相同。）

\begin{enumerate}
\def\labelenumi{\arabic{enumi})}
\setcounter{enumi}{13}
\tightlist
\item
  \begin{enumerate}
  \def\labelenumii{\alph{enumii})}
  \tightlist
  \item
    设 \(G\) 为局部紧群，\(f\) 是 \(G\) 上的一个下半连续正函数，\(\mu\) 是 \(G\) 上的一个正测度。证明函数
  \end{enumerate}
\end{enumerate}

\[
x \mapsto \int^ {*} f (s ^ {- 1} x) d \mu (s) = g (x)
\]

在 G 上是下半连续的（cf.~Ch. IV, §1, No.~1, Th. 1）；为使 \(\mu\) 与 \(f\) 可卷积，必须且只需 \(g\) 对 G 上的某个左 Haar 测度可积。

\begin{enumerate}
\def\labelenumi{\alph{enumi})}
\setcounter{enumi}{1}
\tightlist
\item
  在群 \(\mathbf{G} = \mathbf{R} \times \mathbf{R}\) 上，设 \(\mu\) 是测度 \(\varepsilon_0 \otimes \lambda\) ，其中 \(\lambda\) 是 \(\mathbf{R}\) 上的 Lebesgue 测度；设 \(f(x, y) = (1 - |xy - 2|)^+\) ；证明函数
\end{enumerate}

\[
g (x, y) = \int f (x - s, y - t) d \mu (s, t)
\]

在 G 上处处有限，但既不连续，也不对 G 上的 Lebesgue 测度可积。

¶ 15) 设 G 为局部紧群，β 为 G 上的一个左 Haar 测度；以下按语言上的滥用，把 β-可积的数值函数与它们在 L¹(G,β) 中的类等同起来；对 Lᵖ(G,β) 同样滥用。

\begin{enumerate}
\def\labelenumi{\alph{enumi})}
\item
  设 \(A\) 是 Banach 空间 \(\mathrm{L}^1(\mathrm{G},\beta)\) 的一个连续自同态，使得对每个 \(s \in \mathrm{G}\) 都有 \(A(f*\varepsilon_s) = A(f)*\varepsilon_s\) （对所有 \(f \in \mathrm{L}^1(\mathrm{G},\beta)\) ）。证明：于是对每个函数 \(g \in \mathcal{K}(\mathrm{G})\) 都有 \(A(f*g) = A(f)*g\) （注意 \(s \mapsto g(s)A(f*\varepsilon_s)\) 是 \(\beta\) -可积的映射（\(\mathrm{G}\) 到 \(\mathrm{L}^1(\mathrm{G},\beta)\) 中的））；反之亦然。由此推出：存在唯一的有界测度 \(\mu\) （\(\mathrm{G}\) 上的），使得 \(A(f) = \mu *f\) 对所有 \(f \in \mathrm{L}^1(\mathrm{G},\beta)\) 成立，且 \(\| A\| = \| \mu \|\) 。（用 No.~7 的命题 19 的记号，取 \(A(f_{\mathrm{V}}*g)\) 的极限（关于比 \(\mathfrak{B}\) 的截口滤子更细的一个超滤子），并利用 \(\mathcal{M}^1(\mathrm{G})\) 的单位球对弱拓扑 \(\sigma (\mathcal{M}^1(\mathrm{G}),\overline{\mathcal{K}(\mathrm{G})})\) 的紧性。）
\item
  设 \(\mu\) 是 G 上的一个有界测度；直接证明连续自同态 \(\pmb{\gamma}(\mu): f \mapsto \mu * f\) （\(\mathrm{L}^{\infty}(\mathrm{G},\beta)\) 的）的范数等于 \(\| \mu \|\) 。（化归到 \(\mu\) 具有紧支集且关于 \(|\mu|\) 具有连续密度的情形。）由此重新推出连续自同态 \(\pmb{\gamma}(\mu): f \mapsto \mu * f\) （\(\mathrm{L}^{1}(\mathrm{G},\beta)\) 的）的范数等于 \(\| \mu \|\) 。
\item
  假设 G 是紧的且 \(\mu\) 是正的。证明对 \(1 < p < +\infty\) ，连续自同态 \(\gamma(\mu): f \mapsto \mu * f\) （\(\mathrm{L}^{p}(\mathrm{G}, \beta)\) 的）的范数等于 \(\|\mu\|\) 。
\item
  取 G 为 3 阶循环群。给出 G 上一个测度 \(\mu\) 的例子，使得自同态 \(\pmb{\gamma}(\mu)\) （\(\mathrm{L}^p (\mathrm{G},\beta)\) 的）的范数严格小于 \(\| \mu \|\) （对 \(1 < p < +\infty\) ）。
\end{enumerate}

¶ 16) 记号与约定均是 Exer. 15 的。

\begin{enumerate}
\def\labelenumi{\alph{enumi})}
\tightlist
\item
  证明：为使有界测度 \(\mu\) （\(\mathbf{G}\) 上的）满足 \(\| \mu * f \|_1 = \| f \|_1\) 对所有 \(f \in \mathrm{L}^1(\mathbf{G}, \beta)\) 都成立，必须且只需 \(\mu\) 是范数为 1 的点测度。（利用自同态 \(\gamma(|\mu|)\) （\(\mathrm{L}^1(\mathbf{G}, \beta)\) 的）的范数为 \(\| \mu \|\) 这一事实（Exer. 15），证明必然有：对每个函数 \(f \in \mathcal{K}(\mathbf{G})\) ，
\end{enumerate}

\[
\left| \int f d \mu \right| = \int | f | d | \mu |;
\]

由此先推出 \(\mu = c|\mu|\) ，其中 \(c\) 是模为 1 的常数，再推出 \(\mu\) 是一个点测度。）

\begin{enumerate}
\def\labelenumi{\alph{enumi})}
\setcounter{enumi}{1}
\tightlist
\item
  取 G 为 3 阶循环群。给出 G 上一个测度 \(\mu\) 的例子，它不是点测度，使得 \(\| \mu * f \|_2 = \| f \|_2\) 对 G 上定义的每个数值函数 \(f\) 都成立。
\end{enumerate}

\begin{enumerate}
\def\labelenumi{\arabic{enumi})}
\setcounter{enumi}{16}
\tightlist
\item
  记号与约定均是 Exer. 15 的。
\end{enumerate}

\begin{enumerate}
\def\labelenumi{\alph{enumi})}
\item
  设 \(\mu\) 是 G 上的一个有界测度；为使自同态 \(\pmb{\gamma}(\mu): f \mapsto \mu * f\) （\(\mathrm{L}^p(\mathrm{G}, \beta)\) 的）（ \(1 \leqslant p < +\infty\) ）是满射的，必须且只需存在一个 \(c_p > 0\) ，使得 \(\| \pmb{\mu}^{\vee} * g \|_q \geqslant c_p \| g \|_q\) 对所有 \(g \in \mathrm{L}^q(\mathrm{G}, \beta)\) 都成立（其中 \(\frac{1}{p} + \frac{1}{q} = 1\) ）（cf.~TVS, IV, §4, No.~2, Th. 1 的 Cor. 3 之后的注）。
\item
  若 \(\mu\) 是一个以 \(\beta\) 为基的测度，且 G 不是离散的，证明 \(\gamma(\mu)\) 对 \(1 \leqslant p \leqslant +\infty\) 永远不是满射的（对 \(p < +\infty\) ，利用 \(a\) ，化归到 \(\mu\) 关于 \(\beta\) 的密度属于 \(\mathcal{K}(\mathrm{G})\) 的情形；对 \(p = +\infty\) ，注意对 \(f \in \mathrm{L}^1(\mathrm{G}, \beta)\) 与 \(g \in \mathrm{L}^\infty(\mathrm{G}, \beta)\) ，\(f * g\) 对右一致结构是一致连续的，由此直接论证）。
\item
  若 G 是交换的，\(1 \leqslant p \leqslant 2\) ，且自同态 \(\pmb{\gamma}(\mu)\) （\(\mathrm{L}^p(\mathrm{G},\beta)\) 的）是满射的，证明它是双射的（用正则化证明：若 \(\pmb{\gamma}(\mu)\) 不是单射的，则自同态 \(\pmb{\gamma}(\check{\mu})\) （\(\mathrm{L}^q(\mathrm{G},\beta)\) 的）不是单射的：利用 Ch. IV, §6, No.~5, Prop. 4 的推论）。
\item
  取 \(\mathbf{G} = \mathbf{Z}\) 且 \(\mu = \varepsilon_1 - \varepsilon_0\) ；证明 \(\pmb{\gamma}(\mu)\) 在 \(\mathrm{L}^p (\mathbf{Z},\beta)\) （凡 \(p\neq +\infty\) 者）中都是单射的，但在 \(\mathrm{L}^{\infty}(\mathbf{Z},\beta)\) 中不是，而且对 \(p\) 的任何值都不是满射的。
\end{enumerate}

¶ 18) 记号与约定均是 Exer. 15 的。

\begin{enumerate}
\def\labelenumi{\alph{enumi})}
\item
  设 \(\mu\) 是 \(G\) 上的一个有界测度，其支集至少含两个不同的点。证明存在一个紧集 \(K\) 和一个数 \(k\)（\(0 < k < 1\)），使得 \(\| \varphi_{sK} \cdot \mu \| \leqslant k \| \mu \|\) 对所有 \(s \in G\) 都成立。（若 \(t, t'\) 是 \(\mu\) 的支集里两个不同的点，取 \(K\) 为 \(e\) 的一个充分小的邻域，使 \(sK\) 不能同时与 \(tK\) 和 \(t'K\) 相交。）
\item
  设 \(\mu\) 是一个有界测度 \(\geqslant 0\) （\(G\) 上的），其支集至少含两个不同的点。证明存在一个函数 \(f \in L^{\infty}(G, \beta)\) ，它不等价于任何 \(\geqslant 0\) 的函数，使得 \(\mu * f\) 局部几乎处处 \(\geqslant 0\) （对 \(\beta\) 而言）。（利用 \(a\) ，取 \(f\) 等于 \(-1\) （在 \(K\) 上），等于一个适当的正常数（在 \(\mathbb{C}K\) 上）。）此外，若 \(\mu\) 的支集是紧的，则存在一个具有前面诸性质且具有紧支集的函数 \(f\)。
\item
  设 \(\mu_{1},\mu_{2}\) 是 G 上两个非零的有界正测度，对卷积可交换，使得 \(\mu_{1}\) 的支集 K 是紧的且含有 \(e\) 。令 \(\mu = \mu_{1} + \mu_{2}\) 。设 V 是含有 K 的 \(e\) 的一个对称紧邻域；设 \(g\) 为这样一个函数：它等于 \(\mu * \varphi_{\mathrm{V}}\) （在 \(\mathbb{C}\mathrm{V}\) 上），等于 0（在 V 上）；证明函数 \(\mu * (g - (\mu_1 * \varphi_{\mathrm{V}}))\) 局部几乎处处 \(\geqslant 0\) （在 \(\mathbb{C}(\mathrm{V}^2)\) 中）。另一方面，证明存在一个紧集 H，使函数 \(\mu_{2} * \varphi_{\mathrm{H}}\) 几乎处处 \(\geqslant 0\) （在 \(\mathrm{V}^2\) 中）。
\item
  由 c) 推出：若 \(\mu\) 是 G 上的一个有界正测度，其支集至少含两个不同的点，则在再附加下列假设之一时，存在一个属于每个 \(L^{p}(G,\beta)\) （ \(1 \leqslant p \leqslant +\infty\) ）的函数 f，它不等价于任何 \(\geqslant 0\) 的函数，且使 \(\mu * f\) 局部几乎处处 \(\geqslant 0\) ：\(\alpha\) ) G 是交换的；\(\beta\) ) 存在一点 \(a \in G\) 使 \(\mu(\{a\}) > 0\) 。（把 \(\mu\) 适当地分解为两个可交换的测度 \(\geqslant 0\) 之和。）
\end{enumerate}

¶ 19) 记号与约定均是 Exer. 15 的。

\begin{enumerate}
\def\labelenumi{\alph{enumi})}
\item
  证明：若一个正有界测度 \(\mu\) （\(G\) 上的）使得对每个函数 \(f \in L^{p}(G, \beta)\) （ \(p\) 给定，\(1 \leqslant p < \infty\) ），关系 \(\mu * f \geqslant 0\) 局部几乎处处成立就蕴含 \(f \geqslant 0\) 局部几乎处处成立，则在下列每种情形都可断定 \(\mu\) 是一个点测度：\(\alpha)\) \(G\) 是紧的；\(\beta)\) \(G\) 是离散的；\(\gamma)\) \(G\) 是交换的（利用 Exer. 18) (*)。对 \(p = +\infty\) ，无须对 \(G\) 作补充假设，同样的结论成立。
\item
  设 \(\mu\) 是一个正有界测度（\(G\) 上的），使得：\(1^{\circ} \pmb{\gamma}(\mu)\) 是 \(L^{1}(G, \beta)\) 的一个满射自同态；\(2^{\circ}\) 对每个函数 \(f \in L^{1}(G, \beta)\) ，关系 \(\mu * f \geqslant 0\) 局部几乎处处成立蕴含 \(f \geqslant 0\) 局部几乎处处成立。证明 \(\mu\) 这时是一个点测度。（注意对每个函数 \(g \in L^{\infty}(G, \beta)\) ，关系 \(\stackrel{\vee}{\mu} * g \geqslant 0\) 局部几乎处处成立蕴含 \(g \geqslant 0\) 局部几乎处处成立。）
\end{enumerate}

\begin{enumerate}
\def\labelenumi{\arabic{enumi})}
\setcounter{enumi}{19}
\tightlist
\item
  设 \(G\) 为局部紧群，\(\beta\) 为 \(G\) 上的一个左 Haar 测度。
\end{enumerate}

\begin{enumerate}
\def\labelenumi{\alph{enumi})}
\item
  设 \(f\) 是 \(G\) 上的一个有界数值函数，对左一致结构一致连续。证明对每个有界测度 \(\mu\) （\(G\) 上的），\(\mu * f\) 都对左一致结构一致连续；此外，用 No.~7 的命题 19 的记号，\(f\) （对 \(G\) 中的一致收敛拓扑而言）是函数 \(f_V * f\) 关于 \(\mathfrak{B}\) 的截口滤子的极限。
\item
  取 \(\mathbf{G} = \mathbf{T}\) ；给出一个连续函数 \(h\) （\(\mathbf{T}\) 上的）的例子，它不是 \(f*g\) 的形式，其中 \(f\) 与 \(g\) 属于 \(\mathrm{L}^2 (\mathbf{T},\beta)\) 。（利用 Ch. IV, §6 的 Exers. 15 b) 与 16。）
\end{enumerate}

¶ 21) 设 G 为局部紧群，β 为 G 上的一个左 Haar 测度；把 L¹(G,β) 典范地等同于 M¹(G) 的一个子空间。

\begin{enumerate}
\def\labelenumi{\alph{enumi})}
\item
  用 §3 的 Exer. 11 的记号，设 A 是 \(\mathcal{M}_{\mathrm{III}}\) 的一个相对紧子集，B 是 \(\mathrm{L}^1 (\mathrm{G},\beta)\) 的一个在 \(\mathcal{M}_{\mathrm{IV}}\) 中相对紧的子集。证明 \(\mathrm{A}\times \mathrm{B}\) 上映射 \((\mu ,\nu)\mapsto \mu *\nu\) （\(\mathcal{M}_{\mathrm{III}}\times \mathcal{M}_{\mathrm{IV}}\) 到 \(\mathcal{M}_{\mathrm{IV}}\) 中的）的限制是连续的。（先证明若 A 与 B 是紧的，则 \(\mathrm{A}\times \mathrm{B}\) 在此映射下的像是紧的；利用 §3 的 Exer. 10 以及 Ch. V, §5, Exer. 15 的判别准则 \(\alpha\) )。为再证明 \((\mu ,\nu)\mapsto \mu *\nu\) 连续，利用 §3 的 Exer. 11 c)。）
\item
  设 \(\mathcal{U}_s^\infty (\mathrm{G})\) 是 \(\mathbf{G}\) 上对左一致结构一致连续的有界数值函数的集合。用 \(\mathcal{T}_{\mathrm{II}}\) 记拓扑 \(\sigma (\mathcal{M}^1 (\mathrm{G}),\mathcal{U}_s^\infty (\mathrm{G}))\) ，并用 \(\mathcal{M}_{\mathrm{II}}\) 记空间 \(\mathcal{M}^1 (\mathrm{G})\) 赋以 \(\mathcal{T}_{\mathrm{II}}\) 后所得的空间。取 \(\mathbf{G} = \mathbf{R}\) ，给出测度序列 \((\mu_n)\) （\(\mathbf{G}\) 上的）对 \(\mathcal{T}_{\mathrm{II}}\) 趋于 0 的一个例子，以及函数序列 \((f_n)\) （\(\mathrm{L}^1 (\mathrm{G},\beta)\) 中的）对 \(\mathcal{T}_{\mathrm{IV}}\) （或者等价地，对拓扑 \(\sigma (\mathrm{L}^1 (\mathrm{G},\beta),\mathrm{L}^\infty (\mathrm{G},\beta))\) ）趋于 0 的一个例子，使得序列 \((\mu_n*f_n)\) 对 \(\mathcal{T}_{\mathrm{IV}}\) 不趋于 0（取 \(f_{n}(t) = \sin nt\) 于区间 \([0,\pi ]\) 中，取 \(f_{n}(t) = 0\) 于别处）。
\item
  设 A 是 \(\mathcal{M}_{\mathrm{II}}\) 的一个相对紧子集，B 是 \(\mathrm{L}^1 (\mathrm{G},\beta)\) 的一个对范数 \(\| f\| _1\) 的拓扑相对紧的子集。证明 \(\mathbf{A}\times \mathbf{B}\) 上映射 \((\mu ,f)\mapsto \mu *\textit{f}\) （\(\mathcal{M}_{\mathrm{III}}\times \mathrm{L}^1 (\mathrm{G},\beta)\) 到 \(\mathrm{L}^1 (\mathrm{G},\beta)\) 中的）的限制是连续的（L \(^1 (\mathrm{G},\beta)\) 赋以其赋范空间拓扑）。（化归为证明：对 \(f\in \mathcal{K}(\mathrm{G})\) ，映射 \(\mu \mapsto \mu *\textit{f}\) （\(\mathcal{M}_{\mathrm{II}}\) 到 \(\mathrm{L}^1 (\mathrm{G},\beta)\) 中的）在 A 上连续。利用对 \(g\in \mathrm{L}^{\infty}(\mathrm{G},\beta)\) 与 \(f\in \mathcal{K}(\mathrm{G})\) ，\(g*\textit{f}\) 对左一致结构一致连续这一事实，先证明映射 \(\mu \mapsto \mu *\textit{f}\) （\(\mathcal{M}_{\mathrm{II}}\) 到 \(\mathcal{M}_{\mathrm{IV}}\) 中的）是连续的。然后，利用 Exer. 20 a) ，化归为证明：若 C 是 \(\mathrm{L}^1 (\mathrm{G},\beta)\) 的一个对拓扑 \(\sigma (\mathrm{L}^{1}(\mathrm{G},\beta),\mathrm{L}^{\infty}(\mathrm{G},\beta))\) 紧的子集，且 \(h\in \mathcal{K}(\mathbf{G})\) ，则映射 \(\mu \mapsto \mu *h\) （C 到 \(\mathrm{L}^1 (\mathrm{G},\beta)\) 中的）当 \(\mathrm{L}^1 (\mathrm{G},\beta)\) 赋以其赋范空间拓扑时是连续的。用与 \(a\) ) 中同样的推理证明：为此只需证明 C 在此映射下的像对赋范空间拓扑是紧的。为此，利用 Šmulian 定理、C 是挤压的（§3, Exer. 10）这一事实以及 Lebesgue 定理。）
\item
  设 A 是 \(\mathcal{M}_{\mathrm{II}}\) 的一个含有 0 的相对紧子集，B 是 \(\mathcal{M}_{\mathrm{I}}\) 的一个相对紧子集；证明对每个 \(\nu_0 \in \mathrm{B}\) ，\(\mathrm{A} \times \mathrm{B}\) 上映射 \((\mu, \nu) \mapsto \mu * \nu\) （\(\mathcal{M}_{\mathrm{II}} \times \mathcal{M}_{\mathrm{I}}\) 到 \(\mathcal{M}_{\mathrm{II}}\) 中的）的限制在点 \((0, \nu_0)\) 处连续。（利用 Exer. 20 a) 与 Exer. 21 c)。）
\item
  设 A 是 \(\mathcal{M}^1 (\mathrm{G})\) 的一个有界子集，\(f\) 是 \(\mathrm{L}^1 (\mathrm{G},\beta)\) 中的一个函数；证明映射 \(\mu \mapsto \mu *f\) （\(\mathcal{M}_{\mathrm{III}}\) 到 \(\mathcal{M}_{\mathrm{IV}}\) 中的）在 A 上的限制是连续的。
\item
  设 \(\mathcal{M}_{+}^{1}(G)\) 是 G 上有界正测度的集合。证明：若 \(\mathcal{M}_{+}^{1}(G)\) 的子集 A 对 \(T_{II}\) 是紧的，则它对 \(T_{III}\) 也是紧的。（化归为证明 A 是一个挤压集。用反证法，利用对 \(f \in \mathcal{K}(G)\) ，A 在映射 \(\mu \mapsto \mu * f\) 下的像依 c) 是一个挤压集这一事实。）
\end{enumerate}

\(g)\) 证明对 \(\mathbf{G} = \mathbf{R}\) ，\(\mathcal{M}_{+}^{1}(\mathbf{G})\) 上由 \(\mathcal{T}_{\mathrm{II}}\) 与 \(\mathcal{T}_{\mathrm{III}}\) 诱导的拓扑是不同的。

¶ 22) 记号是 §4 的 Exer. 21 与 §3 的 Exer. 11 的。

\begin{enumerate}
\def\labelenumi{\alph{enumi})}
\item
  设 \((\mu_n)\) 是 \(G\) 上的有界测度序列。假设对每个函数 \(f \in L^1(G, \beta)\) ，序列 \((\mu_n * f)\) 在 \(\mathcal{M}_I\) 中收敛于 0。证明范数序列 \((\|\mu_n\|)\) 是有界的（对映射族 \(f \mapsto \mu_n * f\) （\(L^1(G, \beta)\) 到自身的）利用 Banach-Steinhaus 定理，并利用 Exer. 15）。由此推出序列 \((\mu_n)\) 在 \(\mathcal{M}_I\) 中趋于 0（利用 Exer. 20）。
\item
  设 \((\mu_n)\) 是 \(G\) 上的有界测度序列，使得序列 \((\mu_n * f)\) 对每个函数 \(f \in L^1(G, \beta)\) 淡收敛于 0；证明序列 \((\mu_n)\) 淡收敛于 0（对每个紧子集 \(K\) （\(G\) 的），按 \(a\) ) 的方式论证序列 \((|\mu_n|(K))\) 有界）。给出一个例子（取 \(G = Z\) ），其中序列 \((\|\mu_n\|)\) 不是有界的。
\item
  设 \((\mu_n)\) 是 \(G\) 上的有界测度序列，使得对每个函数 \(f \in L^1(G, \beta)\) ，序列 \((\mu_n * f)\) 在 \(\mathcal{M}_{II}\) 中收敛于 0；证明序列 \((\mu_n)\) 在 \(\mathcal{M}_{II}\) 中收敛于 0（利用 \(a\) )）。
\end{enumerate}

\begin{enumerate}
\def\labelenumi{\arabic{enumi})}
\setcounter{enumi}{22}
\item
  设 \(G\) 为局部紧群，\(\beta\) 为 \(G\) 上的一个左 Haar 测度，\(f\) 与 \(g\) 是两个局部 \(\beta\) -可积的正函数。假设 \(f\) 与 \(g\) 可卷积，且其中一个函数在可数多个紧集之并的余集上为零。证明 \(f * g\) 局部几乎处处等于一个下半连续函数（利用 No.~5 的命题 15）。
\item
  证明 No.~5 的命题 12 与 15 中的不等式不能通过在其右端插入一个常数因子 c \textless{} 1 而加以改进。
\item
  设 G 为局部紧群，\(\beta\) 为 G 上的一个左 Haar 测度，\(\mu\) 为 G 上的一个正测度。若 A 是一个 \(\mu\) -可积集而 B 是 G 中的一个 Borel 集，则函数
\end{enumerate}

\[
u: s \mapsto \mu (\mathrm{A} \cap s \mathrm{B})
\]

在 G 中是 \(\beta\) -可测的；此外，若 \(B^{-1}\) 是 \(\beta\) -可积的，则 \(u\) 也是，且

\[
\int \mu (\mathrm{A} \cap s \mathrm{B}) d \beta (s) = \mu (\mathrm{A}) \beta (\mathrm{B} ^ {- 1}).
\]

给出一个测度 \(\mu\) 的例子，使得 \(s \mapsto \mu(A \cap sB)\) 不连续；若 \(\mu\) 是一个以 \(\beta\) 为基的测度，且 A 是 \(\mu\) -可积的而 B 是一个 Borel 集，则函数 \(s \mapsto \mu(A \cap sB)\) 在 G 上连续。

¶ 26) 设 G 为局部紧群，β 为 G 上的一个左 Haar 测度。为使 L\^{}p(G,β) 的子集 H (1 ≤ p \textless{} +∞) 是相对紧的（对 p 阶平均收敛拓扑），必须且只需下列条件满足：1° H 在 L\^{}p(G,β) 中有界；2° 对每个 ε \textgreater{} 0，存在 G 的紧子集 K，使得对每个函数 f ∈ H 有 \textbar\textbar fφ\_G-K\textbar\textbar\_p ≤ ε；3° 对每个 ε \textgreater{} 0，存在 G 中 e 的邻域 V，使得对所有 f ∈ H 与 s ∈ V 有 \textbar\textbar γ(s)f - f\textbar\textbar\_p ≤ ε。（为证明这些条件是充分的，注意：若 g ∈ K(G) 且 L 是 G 的紧子集，则集合 φ\_L · H 在映射 f ↦ g * f 下的像是 K(G) 的一个等度连续子集。）

¶ 27) 设 G 为局部紧群，β 为 G 上的一个左 Haar 测度。若 G 不约化为 e，则 L¹(G,β) 是一个（对卷积的）具有异于 0 的零因子的代数。为作出 L¹(G,β) 中两个使 f * g = 0 的非零元素 f, g，可按如下方式进行：

\(1^{\circ}\) \(\mathbf{G}\) 含有一个紧子群 \(\mathbf{H}\) （不约化为 \(e\) ），且在其中 \(\Delta(x)=1\) 的情形。取 \(f\) 为一个集合 \(\mathbf{A}\) 的特征函数，取 \(g\) 为两个特征函数之差 \(\varphi_{s\mathbf{B}}-\varphi_{\mathbf{B}}\) ，其中 \(\mathbf{A}\) 与 \(\mathbf{B}\) 适当选取。

\(2^{\circ}\) \(\mathbf{G} = \mathbf{Z}\) 的情形。证明这时可取

\[
f (n) = \frac {1}{2 n - 1} - \frac {1}{2 n + 1}
\]

对所有 \(n \in \mathbf{Z}\) ，以及 \(g(n) = f(-n)\) 。

\(3^{\circ}\) 一般情形。先证明 G 中存在一个 \(a \neq e\) 使 \(\Delta(a) = 1\) 。于是 \(a\) 生成的子群在 G 中的闭包 H 或者是一个紧子群，或者是一个同构于 Z 的子群 (GT, V, §1, Exer. 2)。第一种情形用 \(1^{\circ}\) 的结果；第二种情形取

\[
f (t) = \sum_ {n = - \infty} ^ {+ \infty} \alpha_ {n} \varphi_ {\mathrm{U} a ^ {- n}} (t), \quad g (t) = \sum_ {n = - \infty} ^ {+ \infty} \beta_ {n} \varphi_ {a ^ {n} \mathrm{U}} (t),
\]

适当选取 U 与序列 \((\alpha_{n})\) 、\((\beta_{n})\) （借助 \(2^{\circ}\) ）。

\begin{enumerate}
\def\labelenumi{\arabic{enumi})}
\setcounter{enumi}{27}
\tightlist
\item
  设 \(G_1, G_2\) 为两个局部紧群，\(\beta_1, \beta_2\) 为 \(G_1, G_2\) 上的左 Haar 测度，\(A_1\) （相应地 \(A_2\) ）为（ \(R\) 上的）拓扑代数 \(L^1(G_1, \beta_1)\) （相应地 \(L^1(G_2, \beta_2)\) ）。设 \(T\) 是 \(A_1\) 到 \(A_2\) 上的一个代数同构，使得关系 \(f \geqslant 0\) 几乎处处成立等价于 \(T(f) \geqslant 0\) 几乎处处成立。
\end{enumerate}

\begin{enumerate}
\def\labelenumi{\alph{enumi})}
\item
  证明 \(\mathbf{T}\) 是拓扑代数的一个同构。（先注意：若一个递减序列 \((f_n)\) （由 \(\mathbf{A}_1\) 的元素组成）几乎处处趋于 0，则诸 \(\mathrm{T}(f_n)\) 构成的序列也几乎处处趋于 0；由此推出：对每个序列 \((f_n)\) （\(\mathbf{A}_1\) 中有上界的），有 \(\mathrm{T}(\sup (f_n)) = \sup (\mathrm{T}(f_n))\) ，并借助 Fatou 引理得出结论。）
\item
  证明 T 可以且只可以用一种方式扩张为拓扑代数 \(\mathcal{M}^{1}(G_{1})\) 到拓扑代数 \(\mathcal{M}^{1}(G_{2})\) 上的同构，且存在 \(G_{1}\) 到 \(G_{2}\) 上的一个拓扑同构 u 与一个连续同态 \(\chi\) （\(G_{1}\) 到 \(R_{+}^{*}\) 中的），使得 \(\mathrm{T}(\varepsilon_{s}) = \chi(s)\varepsilon_{u(s)}\) （利用 Exer. 15 a) 与 Exer. 19 b)；为证明 u 与 \(\chi\) 连续，注意 T 定义了代数 \(\mathcal{L}(A_{1})\) （拓扑向量空间 \(A_{1}\) 的连续自同态所成的）到类似的代数 \(\mathcal{L}(A_{2})\) 上的一个同构，且这个同构对逐点收敛拓扑是双连续的；另一方面，注意 \(s \mapsto \delta(s^{-1})\) 是 \(G_{1}\) 到 \(\mathcal{L}(A_{1})\) 的一个乘法子群上的一个同构。
\item
  由 \(b)\) 推出：对每个 \(f \in \mathbf{A}_1\) ，
\end{enumerate}

\[
(T (f)) (t) = \chi (u ^ {- 1} (t)) f (u ^ {- 1} (t))
\]

对所有 \(t \in \mathbf{G}_2\) 成立。

回忆 (A, III, §2, Exer. 9)：存在有限群 \(\mathbf{G}_1, \mathbf{G}_2\) ，使得代数 \(\mathrm{L}^1(\mathrm{G}_1, \beta_1)\) 与 \(\mathrm{L}^1(\mathrm{G}_2, \beta_2)\) 同构而 \(\mathbf{G}_1\) 与 \(\mathbf{G}_2\) 并不同构。

